% !TEX TS-program = pdflatex
% !TeX program = pdflatex
% !TEX encoding = UTF-8
% !TEX spellcheck = en_US

\documentclass[11pt, a4paper]{article}
%\usepackage{fullpage}
\usepackage[left=1cm,right=1cm,top=1cm,bottom=2cm]{geometry}
\usepackage[fleqn]{amsmath}
\usepackage{amssymb}
%\usepackage{indentfirst}
\usepackage[T1]{fontenc}
\usepackage[utf8]{inputenc}
\usepackage[french,english]{babel}
\usepackage{txfonts} 
\usepackage[]{graphicx}
\usepackage{multirow}
\usepackage{hyperref}
\usepackage{parskip}
\usepackage{multicol}
\usepackage{wrapfig}

\usepackage{turnstile}%Induction symbole

\usepackage{tikz}
\usetikzlibrary{arrows, automata}
\usetikzlibrary{decorations.pathmorphing}

\renewcommand{\baselinestretch}{1}

\setlength{\parindent}{0pt}


\begin{document}

%\selectlanguage {french}
%\pagestyle{empty} 

\noindent
\begin{tabular}{ll}
	\multirow{3}{*}{\includegraphics[width=1.5cm]{../../extra/logo/esi.nlp.pdf}} & \'Ecole national Supérieure d'Informatique\\
	& 2\textsuperscript{nd} year second cycle (2CSSID)\\
	& NLP: Natural Language Processing (2022-2023)
\end{tabular}\\[.25cm]
\noindent\rule{\textwidth}{2pt}\\[-0.5cm]
\begin{center}
	{\LARGE \textbf{Workshop 01: Sentiment analysis}}
	\begin{flushright}
		Abdelkrime Aries
	\end{flushright}
\end{center}\vspace{-0.5cm}
\noindent\rule{\textwidth}{2pt}

This is a guided ungraded lab. 
Its purpose of workshops is to understand how to use some high level API's to solve a problem based on some demos about these API's.
The problem we want to solve in this workshop is a well known one: text-based sentiment analysis.
Our aim is to understand how to manipulate different representations of words/sentences.


\section*{General information}

\subsection*{Data}

Financial sentiment analysis:
\begin{itemize}
	\item URL: \url{https://www.kaggle.com/datasets/sbhatti/financial-sentiment-analysis}
	\item Description: Some texts about finance with their sentiment orientation (positive, negative, neutral).
	It contains 860 negative samples, 1852 positive ones, and 3130 neutral ones.
\end{itemize}

Sentiment Lexicons for Text Mining:
\begin{itemize}
	\item URL: \url{https://www.kaggle.com/datasets/ekrembayar/sentiment-lexicons-for-text-mining}
	\item Description: A lexicon of English words with their sentiment orientation.
\end{itemize}

\subsection*{Tools}

NLTK, Scikit-learn, numpy, pandas, tensorflow, gensim

%\begin{itemize}
%	\item NLTK
%	\item Scikit-learn
%	\item numpy
%	\item pandas
%	\item tensorflow
%	\item gensim
%\end{itemize}

\subsection*{Demos}

\begin{itemize}
	\item NLTK: \url{https://github.com/projeduc/ESI_TALN/blob/master/tuto/CH06/encoding_python_NLTK.ipynb}
	\item Gensim: \url{https://github.com/projeduc/ESI_TALN/blob/master/tuto/CH06/encoding_python_gensim.ipynb}
	\item Scikit-learn: \url{https://github.com/projeduc/ESI_TALN/blob/master/tuto/CH06/encoding_python_sklearn.ipynb}
	\item Tensorflow Auto-encoder: \url{https://github.com/projeduc/ESI_ML/blob/main/demos/NN/TF_Autoencoder.ipynb} (can be used for MLP)
	\item Tensorflow CNN: \url{https://github.com/projeduc/ESI_ML/blob/main/demos/NN/TF_CNN.ipynb} (CNN for images are similar to those for text)
\end{itemize}


\section{Data preparation}

\begin{itemize}
	\item Import the sentiment lexicon "bing.csv" and transform it into a Python dictionary.
	\item Transform positive label into +1 and negative one into -1.
	\item Import the financial sentiment analysis dataset.
	\item Split this last dataset into train and test using \textbf{train\_test\_split}; the test size is 30\% using random\_state=0 to get a stable split over executions.
	\item Train a \textbf{LabelBinarizer} on training labels and transform them into OneHot representation along with test data.
\end{itemize}

\subsection{TF encoding}

\begin{itemize}
	\item Train a \textbf{CountVectorizer} on the training datasaet, vectorize it along with the test dataset. 
	\item Define a function \textbf{tokenstem} which takes a sentence and tokenize it using \textbf{tokenizer} (defined in the starter project), then each token is stemmed using \textbf{nltk.stem.porter.PorterStemmer}, and last the result is a list of words which do not belong to \textbf{nltk.corpus.stopwords.words('english')}.
	\item Train another \textbf{CountVectorizer} using  \textbf{tokenstem} as \textbf{tokenizer}. 
	\item Train a third one using  \textbf{tokenstem} as \textbf{tokenizer} and limiting the vector's size \textbf{max\_features = 3000}.
\end{itemize}

\subsection{Embedding encoding}

\begin{itemize}
	\item Using the previous function \textbf{tokenizer} which splits a sentence into words, create a list of sentences each is a list of words.	
	\item Train a \textbf{Word2Vec} of \textbf{Gensim} using a window of \textbf{3} and minimum count of \textbf{1}.
	The vector must have a size of \textbf{5}.
	\item Train a \textbf{Fasttext} of \textbf{Gensim} using the same parameters. Fasttext is a character-based word embedding.
	\item When encoding a word, it may be out of vocabulary. So, it will not have a code and we will get an error \textbf{KeyError}.
	In this case, for simplification, we consider the code as a vector of zeros.
	\item Encode each sentence using the centroid of its words (in this case \textbf{size = 5}).
	\item Encode each sentence using the consecutive codes of its words. 
	The maximum sentence size is \textbf{20} so their concatenation will have \textbf{size = 20 * 5 = 100}. 
	If the sentence is longer, it must be truncated.
	If it is less, you must add some padding (a vector of zeros).
	\item Encode each sentence as a matrix of the codes of its words.
	The maximum sentence size is \textbf{20} so the matrix will have \textbf{size = 20X5}.
	If the sentence is longer, it must be truncated.
	If it is less, you must add some padding (a vector of zeros).
\end{itemize}

\section{Training}

\subsection{TF-based sentiment analysis}

\begin{itemize}
	\item Train three \textbf{MLP} models, each for each TF representation. 
	The models must have the same architecture (the inputs are different, so the number of the first hidden layer's parameters).
	You have the choice of parameters: hidden layers, activation functions, etc.
	\item Train three \textbf{Multinomial Naive Bayes} models, each for each TF representation.
\end{itemize}

\subsection{Vector-based sentiment analysis}

\begin{itemize}
	\item Train 4 \textbf{MLP} models for Wor2Vec and Fasttext embeddings with vector representations 5 and 100. 
	\item Train 2 \textbf{CNN} models for Wor2Vec and Fasttext embeddings with matrix representation.
\end{itemize}

\subsection{Implicit Embedding sentiment analysis}

\begin{itemize}
	\item Design a system for sentiment analysis with Embedding layer where the output is 2 floats. 
	\item Train this model and plot the distribution of texts according to their embedding.
\end{itemize}

\subsection{Rule-based sentiment analysis}

\begin{itemize}
	\item Design a system which calculates the orientation of a sentence $S$ as the average of its words orientations. 
	\[Orientation(S) = \frac{\sum_i Orientation(S_i)}{|S|}\]
	\item We use the sentiment lexicon "bing" to get the orientation. If the word does not exist, its orientation is 0.
	\item Using the training dataset, you must estimate the two hyper-parameters $\alpha$ and $\beta$ so that:
	\[\forall S \in Train\_dataset, Sentiment(S) = 
	\begin{cases}
		Postive & if\,  Orientation(S) \ge \alpha \\
		Negative & if\,  Orientation(S) \le \beta \\
		Neutral & otherwise  \\
	\end{cases}
	\]
\end{itemize}

\section{Testing}

\begin{itemize}
	\item Predict all classes of test dataset based on all models.
	\item For each one, show the classification report.
	\item Complete Table \ref{tab:comp} with the result you get and discuss these results.
\end{itemize}

\begin{table}[hpt]
	\centering
	\begin{tabular}{lllllllll}
		\hline\hline
		Rep. & Arch. & $ P_+ $ & $ R_+ $ & $ P_- $ & $ R_- $ & $ P_0 $ & $ R_0 $ & Accuracy \\
		\hline
		
		\multirow{2}{*}{TF} & MLP &  &  &  &  &  &  &  \\
		 & MNB &  &  &  &  &  &  &  \\
		
		\multirow{2}{*}{TF-stem} & MLP &  &  &  &  &  &  &  \\
		 & MNB &  &  &  &  &  &  &  \\
		
		\multirow{2}{*}{TF-stem-limit} & MLP &  &  &  &  &  &  &  \\
		 & MNB &  &  &  &  &  &  &  \\
		
		Word2Vec-centroid & MLP  &  &  &  &  &  &  &  \\
		Word2ec-concat & MLP  &  &  &  &  &  &  &  \\
		Fasttext-centroid & MLP  &  &  &  &  &  &  &  \\
		Fasttext-concat & MLP  &  &  &  &  &  &  &  \\
		Word2Vec-matrix & CNN  &  &  &  &  &  &  &  \\
		Fasttext-matrix & CNN  &  &  &  &  &  &  &  \\
		Implicit-emb & MLP  &  &  &  &  &  &  &  \\
		Rule-based & /  &  &  &  &  &  &  &  \\
		\hline\hline
	\end{tabular}
	\caption{Comparison between different architectures.}
	\label{tab:comp}
\end{table}




\end{document}
